\chapter{Metodologia}
\label{cap:metodologia}

O percurso metodológico adotado na pesquisa articula três dimensões complementares: uma Revisão Sistemática da Literatura (RSL), uma campanha experimental controlada e a proposta de um modelo quantitativo de apoio à decisão arquitetural. Essa organização permite combinar evidências secundárias, obtidas a partir da literatura científica, com evidências primárias, obtidas por meio de experimentação, buscando fundamentar a tomada de decisão arquitetural em sistemas intensivos em dados.

A estratégia metodológica adotada aproxima-se da perspectiva da Engenharia de Software Baseada em Evidências, na qual decisões, práticas e recomendações devem ser apoiadas por evidências sistematicamente coletadas, avaliadas e interpretadas \cite{kitchenham2015evidence}. Nesse sentido, a Revisão Sistemática da Literatura contribui para organizar o conhecimento existente, enquanto a campanha experimental permite produzir evidências empíricas próprias sobre decisões arquiteturais específicas. A integração dessas etapas fornece a base para a proposta de um modelo quantitativo de apoio à decisão.

\section{Caracterização da Pesquisa}
\label{sec:caracterizacao-pesquisa}

Esta pesquisa pode ser caracterizada como aplicada, uma vez que busca produzir conhecimento voltado ao apoio da tomada de decisão arquitetural em sistemas intensivos em dados. O objetivo não é apenas compreender um fenômeno técnico, mas também propor uma forma de apoiar escolhas arquiteturais relacionadas ao processamento, movimentação e análise de grandes volumes de dados.

Quanto aos objetivos, a pesquisa possui caráter exploratório e explicativo. É exploratória porque investiga relações entre características dos dados, estratégias de processamento e métricas de desempenho ainda pouco sistematizadas na literatura. Também é explicativa porque busca compreender como determinadas decisões arquiteturais influenciam o comportamento observado em ambientes de processamento intensivo de dados, especialmente em cenários que envolvem processamento em lote e processamento em fluxo.

Quanto à abordagem, a pesquisa combina métodos qualitativos e quantitativos. A dimensão qualitativa está presente na Revisão Sistemática da Literatura, por meio da identificação, categorização e interpretação das técnicas arquiteturais, métricas e lacunas presentes nos estudos primários. A dimensão quantitativa aparece tanto na síntese exploratória dos efeitos reportados na literatura quanto na análise estatística dos resultados experimentais.

Do ponto de vista dos procedimentos técnicos, a pesquisa envolve revisão sistemática, análise documental, experimentação controlada e modelagem quantitativa. A revisão sistemática foi utilizada para consolidar evidências existentes; a experimentação foi conduzida para avaliar empiricamente decisões arquiteturais específicas; e a modelagem quantitativa será utilizada para propor uma abordagem de apoio à decisão baseada em características observáveis dos dados e do processamento.

\section{Desenho Metodológico}
\label{sec:desenho-metodologico}

A pesquisa foi estruturada em três etapas principais. A primeira etapa corresponde à Revisão Sistemática da Literatura, cujo objetivo foi identificar técnicas, abordagens, métricas e lacunas relacionadas ao \textit{design} de sistemas intensivos em dados. Essa etapa fornece a base conceitual e empírica para compreender o estado da arte e delimitar o problema investigado. A condução da revisão seguiu diretrizes consolidadas para revisões sistemáticas em Engenharia de Software \cite{kitchenham2007guidelines,kitchtchenham2009,kitchenham2015evidence} e utilizou o protocolo PRISMA 2020 para relatar o processo de identificação, triagem, elegibilidade e inclusão dos estudos \cite{page2021prisma}.

A segunda etapa corresponde à campanha experimental. Essa etapa foi planejada a partir das lacunas identificadas na literatura e teve como objetivo avaliar empiricamente o comportamento de estratégias de processamento \textit{batch} e \textit{streaming} em diferentes condições operacionais. O desenho experimental foi orientado por princípios de experimentação em Engenharia de Software, considerando definição de objetivos, planejamento dos fatores experimentais, identificação das variáveis, execução controlada, coleta sistemática de métricas e análise estatística dos resultados \cite{wohlin2012experimentation}. O experimento utilizou tecnologias amplamente adotadas em \textit{pipelines} de dados, incluindo Apache Spark, Spark Structured Streaming e Apache Kafka \cite{zaharia2016apache,armbrust2018structured,kreps2011kafka}.

A terceira etapa consiste na proposta de um modelo quantitativo de apoio à decisão arquitetural. Essa etapa busca investigar como evidências da literatura, resultados experimentais e características informacionais dos dados podem ser integrados em uma abordagem capaz de apoiar a comparação entre alternativas arquiteturais. Para isso, serão considerados conceitos da Teoria da Informação, especialmente entropia e redundância, como possíveis indicadores para caracterizar dados e relacioná-los a métricas observáveis de desempenho \cite{shannon1948mathematical}.

A Figura~\ref{fig:desenho-metodologico} apresenta uma visão geral do desenho metodológico da pesquisa.

\begin{figure}[htbp]
\centering
\begin{tikzpicture}[
node distance=1.2cm, 
etapa/.style={
rectangle,
rounded corners,
draw,
align=center,
text width=4.4cm,
minimum height=1.2cm,
font=\small
},
produto/.style={
rectangle,
draw,
align=center,
text width=5cm,
minimum height=1.2cm,
font=\small
},
seta/.style={-{Latex[length=3mm]}, thick}
]

\node[etapa] (rsl) {Etapa 1 \\ Revisão Sistemática da Literatura};
\node[produto, right=1.2cm of rsl] (rslout) {Técnicas, métricas, evidências e lacunas};

\node[etapa, below=of rsl] (exp) {Etapa 2 \\ Campanha Experimental};
\node[produto, right=1.2cm of exp] (expout) {Resultados empíricos sobre \textit{batch} e \textit{streaming}};

\node[etapa, below=of exp] (modelo) {Etapa 3 \\ Modelagem Quantitativa};
\node[produto, right=1.2cm of modelo] (modeloout) {Apoio à decisão arquitetural};


\draw[seta] (rsl) -- (rslout);
\draw[seta] (exp) -- (expout);
\draw[seta] (modelo) -- (modeloout);

\draw[seta] (rslout.south) -- ++(0,-0.3cm) -| (exp.north);
\draw[seta] (expout.south) -- ++(0,-0.3cm) -| (modelo.north);

\end{tikzpicture}
\caption{Desenho metodológico da pesquisa}
\label{fig:desenho-metodologico}
\end{figure}

A Revisão Sistemática da Literatura será detalhada no Capítulo~\ref{cap:rsl}, enquanto a campanha experimental será apresentada no Capítulo~\ref{cap:campanha-experimental}. Neste capítulo, o objetivo é apresentar a lógica metodológica geral que conecta essas etapas e justifica sua integração.

\section{Procedimentos de Análise}
\label{sec:procedimentos-analise}

A análise da Revisão Sistemática da Literatura foi conduzida em duas dimensões. A primeira dimensão foi qualitativa, envolvendo a leitura, classificação e categorização dos estudos primários. Essa análise permitiu identificar técnicas de \textit{design}, tipos de otimização, métricas utilizadas, estilos arquiteturais recorrentes e lacunas metodológicas. A segunda dimensão foi quantitativa e exploratória, utilizando o \textit{Log Response Ratio} (LRR) como medida normalizada de efeito. Essa escolha permitiu comparar estudos que utilizaram diferentes métricas de desempenho, como tempo de execução, latência, \textit{throughput} e \textit{speedup}.

A análise da campanha experimental foi conduzida a partir das métricas coletadas durante as execuções. Entre as métricas consideradas estão tempo total de execução, \textit{throughput}, latência, uso de CPU, consumo de memória, \textit{backpressure}, taxa de processamento e métricas de entrada e saída dos contêineres. Essas métricas permitiram avaliar o comportamento dos paradigmas \textit{batch} e \textit{streaming} sob diferentes volumes de dados, taxas de eventos, números de \textit{workers} e intervalos de \textit{trigger}.

Para verificar se as diferenças observadas entre cenários eram estatisticamente significativas, foram utilizados procedimentos estatísticos não paramétricos. Inicialmente, foi aplicado o teste de Shapiro-Wilk para avaliar a normalidade das distribuições das métricas coletadas \cite{shapiro1965analysis}. Considerando que distribuições de desempenho em sistemas computacionais nem sempre podem ser assumidas como normalmente distribuídas, foram utilizados testes não paramétricos nas comparações entre cenários.

Para comparações entre dois grupos independentes, foi aplicado o teste de Mann-Whitney U \cite{mann1947test}. Para comparações entre múltiplos grupos, foi utilizado o teste de Kruskal-Wallis \cite{kruskal1952use}. Quando identificadas diferenças significativas entre múltiplos cenários, foram aplicados testes \textit{pós-hoc} de Dunn com correção de Bonferroni \cite{dunn1964multiple}. Além dos testes de significância, foi utilizado o tamanho de efeito de Cliff's Delta para avaliar a magnitude das diferenças entre distribuições \cite{cliff1993dominance}.

A etapa de modelagem quantitativa será conduzida a partir da integração entre os achados da literatura, os resultados experimentais e características informacionais dos dados. A proposta não busca substituir a experimentação, mas oferecer uma forma de apoio preliminar à análise de alternativas arquiteturais. A intenção é investigar se características como volume, taxa de eventos, entropia, redundância e variabilidade podem ser relacionadas a métricas de desempenho, contribuindo para uma tomada de decisão arquitetural mais fundamentada.

\section{Validade da Pesquisa}
\label{sec:validade-pesquisa}

A validade da pesquisa será analisada considerando possíveis ameaças à validade interna, externa, de construção e de conclusão. Essa organização é comum em estudos experimentais em Engenharia de Software e contribui para tornar explícitas as limitações do desenho experimental, da coleta de dados e da interpretação dos resultados \cite{wohlin2012experimentation}.

A validade interna está relacionada à capacidade de atribuir os efeitos observados às variáveis controladas no experimento. Para mitigar ameaças internas, os cenários foram definidos previamente, as execuções foram automatizadas e as métricas foram coletadas de forma sistemática. Além disso, cada cenário foi executado repetidas vezes, reduzindo a influência de variações ocasionais do ambiente.

A validade externa refere-se à possibilidade de generalização dos resultados para outros contextos. Como o experimento utiliza um ambiente específico, baseado em Spark, Kafka e um \textit{dataset} definido, os resultados devem ser interpretados dentro desse contexto. A ampliação para outros \textit{datasets}, tecnologias e ambientes de execução será indicada como possibilidade de trabalhos futuros.

A validade de construção está relacionada à adequação das métricas utilizadas para representar os conceitos investigados. Para mitigar essa ameaça, foram utilizadas métricas amplamente adotadas em avaliações de desempenho, como \textit{throughput}, latência, uso de recursos, \textit{backpressure} e escalabilidade.

A validade de conclusão diz respeito à adequação dos procedimentos estatísticos utilizados para interpretar os resultados. Para reduzir essa ameaça, foram aplicados testes de normalidade, testes não paramétricos, pós-testes e medidas de tamanho de efeito, evitando conclusões baseadas apenas em médias ou observações visuais.

Dessa forma, a metodologia proposta permite que a pesquisa avance de uma compreensão ampla da literatura para uma avaliação empírica específica e, posteriormente, para uma proposta de modelagem quantitativa voltada ao apoio à decisão arquitetural em sistemas intensivos em dados.
